% introduction_coherence.tex
\subsection{Semantic--Behavioral Coherence as a Nomological Relation}

Because both structures are factor solutions over the same items, their comparison has a nomological home: \citet{cronbach-meehl-1955} held that qualitatively different operations measure the same thing when the network ties them to the same construct, and agreement between meaning-based and response-based solutions is just such convergence---disjoint data sources, one needing no respondents, one partition of the items. We call such agreement \emph{semantic--behavioral coherence}, an estimable scale property, not a validity requirement: divergence is expected, and diagnostic, under at least three conditions. First, formative indicator sets: \citet{bollen-lennox-1991} show that the measurement model leaves causal-indicator correlations unconstrained, so semantically related items may barely cluster and an item census, not a sample, is required. Second, causal-system constructs: \citet{borsboom-2008} argues that some clinical constructs are networks of mutually causal symptoms---panic attacks breed worry, worry avoidance---that constitute the disorder rather than measure a common cause, so behavioral covariance tracks the causal chain and can outrun semantic kinship. Third, construct-irrelevant variance: method factors and wording artifacts add behavioral covariance with no semantic counterpart \citep{messick-1995}. There is quantitative precedent: in 39{,}775 responses to the 42-item Depression Anxiety Stress Scales, \citet{stanghellini-etal-2024} found non-random alignment between item-text semantic communities (syntactic and synonym networks) and exploratory-graph-analysis factors, quantified as ``semantic loadings''---Jaccard overlaps---and proposed as a lens on content and discriminant validity. We draw the same comparison with embedding-based semantic structure and explicit congruence statistics, promoting it from diagnostic device to named nomological relation.
